\documentclass[11pt]{article}
\usepackage[a4paper,margin=2.5cm]{geometry}
\usepackage{graphicx}
\usepackage{hyperref}
\usepackage{enumitem}
\usepackage{listings}
\usepackage{xcolor}
\usepackage{amsmath}
\usepackage{amssymb} % or \usepackage{amsfonts}

\title{2dCFT Numerics: Code Structure and Conventions}


\lstset{
basicstyle=\ttfamily\small,
keywordstyle=\color{blue},
commentstyle=\color{gray},
stringstyle=\color{teal},
showstringspaces=false,
breaklines=true,
frame=single
              }

\begin{document}
\maketitle

\section{Purpose}
The purpose of the code structure document is to provide a clear and organized overview of the codebase. We will check the Pentagon equation as well as the idempotency relation for the Fusion kernels. We will first review the spacelike results at rational central charge. Then, we will move on to the timelike results and check the Pentagon. It will amount to use the Virasoro Wick Rotation. We first recall the form of the Pentagon equation and the idempotency relation. The pentagon equation reads
\begin{equation}
   \int_{i\mathbb{R}} \frac{dP_s}{i} \, \mathbb{F}^{(b)}_{P_q, P_4} 
   \begin{bmatrix}
    P_s & P_3 \\
    P_u & P_5
\end{bmatrix}
\mathbb{F}^{(b)}_{P_o, P_s} 
   \begin{bmatrix}
    P_1 & P_2 \\
    P_u & P_q
\end{bmatrix}
\mathbb{F}^{(b)}_{P_s, P_t} 
   \begin{bmatrix}
    P_2 & P_3 \\
    P_1 & P_4
\end{bmatrix}
= \mathbb{F}^{(b)}_{P_o, P_4} 
   \begin{bmatrix}
    P_1 & P_t \\
    P_u & P_5
\end{bmatrix}
\mathbb{F}^{(b)}_{P_q, P_t} 
   \begin{bmatrix}
    P_2 & P_3 \\
    P_o & P_5
\end{bmatrix}
\end{equation}
Ioannis advised me to switch variables $P_s=ip_s$ so that making this change of variable the Pentagon equation reads
\begin{equation}\label{eq:Pentagoninlittlep}
    \begin{aligned}
        &\int_{\mathbb{R}} dp_s \, \mathbb{F}^{(b)}_{P_q, P_4} 
    \begin{bmatrix}
     ip_s & P_3 \\
     P_u & P_5
     \end{bmatrix}
\mathbb{F}^{(b)}_{P_o, ip_s} 
    \begin{bmatrix}
     P_1 & P_2 \\
     P_u & P_q
     \end{bmatrix}
\mathbb{F}^{(b)}_{ip_s, P_t} 
    \begin{bmatrix}
     P_2 & P_3 \\
     P_1 & P_4
     \end{bmatrix}\\
&= \mathbb{F}^{(b)}_{P_o, P_4} 
    \begin{bmatrix}
     P_1 & P_t \\
     P_u & P_5
     \end{bmatrix}
\mathbb{F}^{(b)}_{P_q, P_t} 
    \begin{bmatrix}
     P_2 & P_3 \\
     P_o & P_5
     \end{bmatrix}
    \end{aligned}
\end{equation}
For spacelike Liouville, the fusion kernels on the real line $p\in \mathbb{R}$, \textit{i.e} $P\in i\mathbb{R}$ are real. In this expression, we have unknowns $P_o,P_q, P_1, P_2, P_3, P_4, P_5, P_t$ as well as $b$ parametrized as $b=\sqrt{\frac
{m}{m}}$ (we also define $s=\sqrt{mn}$). We integrate on $p_s$. We will check it for imaginary values of the $P$'s. The idempotency relation reads
\begin{equation}
    \int_{i\mathbb{R}} \frac{dP}{i} \mathbb{F}^{(b)}_{P_s, P} \begin{bmatrix}
     P_2 & P_3 \\
     P_1 & P_4
    \end{bmatrix}
    \mathbb{F}^{(b)}_{P, P_s'} \begin{bmatrix}
     P_3 & P_4 \\
     P_2 & P_1
    \end{bmatrix}
    = \frac{1}{2} \left(\delta(P_s - P_s')+\delta(P_s+P_s')\right)\,.
\end{equation}


\section{Expressions}
\subsection{Special functions}
We first need the expression of the special functions $G_{m,n}$ and $\tilde{G}_{m,n}$. The easiest one is $G_{m,n}$ which is given by the following product
\begin{equation}\label{eq:Gmn}
    G_{m,n}(x)=\prod_{k=0}^{m-1}\prod_{l=0}^{n-1}G\left(x+\frac{k}{m}+\frac{l}{n}\right)\,,
\end{equation}
where $G$ is the Barnes G-function. It should satisfy (consistency check)
\begin{equation}\label{eq:CheckGmn}
    \frac{G_{m,n}(x+1/n)}{G_{m,n}(x)}=m^{1/2 - m x}\,(2\pi)^{(m-1)/2}\Gamma(m x)\,.
\end{equation}
The expression of $\tilde{G}_{m,n}$. We first define $\tilde{G}$ as 
\begin{equation}\label{eq:Gtilde}
    \tilde{G}(x)=\frac{G(1+x)}{G(1-x)}\,,
\end{equation}
and then 
\begin{equation}\label{eq:Gtildemn}
    \begin{aligned}
        \tilde{G}_{m,n}(x)=\prod_{k=0}^{m-1}\prod_{l=0}^{n-1}\tilde{G}\left(x-x_{(k,l)}^{m,n}\right)\,,\\
        x_{(k,l)}^{m,n}=1-\left(\frac{k+1/2}{m}+\frac{l+1/2}{n}\right)\,.
    \end{aligned}
\end{equation}
Finally, we can express the double gamma function 
\begin{equation}
   \Gamma_b(x)=\gamma_{m,n}\frac{(mn)^{\frac{x}{4}(Q-x)}(2\pi)^{\frac{x\sqrt{mn}}{2}-\frac{m+n}{4}}}{G_{m,n}(\frac{x}{\sqrt{mn}})}\,,
\end{equation}
with $\gamma_{m,n}=(mn)^{-\frac{Q^2}{16}}G_{m,n}\left(\frac{1}{2m}+\frac{1}{2n}\right)$. There are few properties of these special functions that we can check numerically(consistency checks):
\begin{itemize}[leftmargin=*]
    \item Complex conjugation:
    \begin{equation}\label{eq:ComplexConjugation}
        \overline{\tilde{G}_{m,n}(x)}=\tilde{G}_{m,n}(\overline{x})\,.
    \end{equation}
    \item Inversion relations:
    \begin{equation}\label{eq:InversionRelations}
        \tilde{G}_{m,n}(x)^{-1}=\tilde{G}_{m,n}(-x)\,.
    \end{equation}
    \item Identity limit:
    \begin{equation}
        \Gamma_b\left(\frac{Q}{2}\right)=1\,.
    \end{equation}
    \item Shift relations:
    \begin{equation}
        \frac{\Gamma_b(x+b^{\pm 1})}{\Gamma_b(x)}=\frac{\sqrt{2\pi}\,b^{\pm b^{\pm 1} x \mp \frac{1}{2}}}{\Gamma(b^{\pm 1} x)}\,.
    \end{equation}
    
\end{itemize}

\subsection{CFT data}
\paragraph{Spacelike case} We choose the following normalization scheme for the 2 point and the 3 point structure constants.

\begin{equation}
    \begin{aligned}
        \langle V_{P_1}(0) V_{P_2}(1)\rangle&=\frac{1}{\rho^{(b)}_0(P_1)}\left(\delta(P_1-P_2)+\delta(P_1+P_2)\right)\,,\\
        \langle V_{P_1}(0) V_{P_2}(1) V_{P_3}(\infty)\rangle&=C_0^{(b)}(P_1,P_2,P_3)\,,
    \end{aligned}
\end{equation}
where 
\begin{equation}
    \begin{aligned}
        \rho^{(b)}_0(P)&=-4\sqrt{2}\sin(2\pi b P)\sin(2\pi P/b)\,,\\
        C_0^{(b)}(P_1,P_2,P_3)&= \frac{\Gamma_b(2Q)\Gamma_b(\frac{Q}{2}\pm P_1 \pm P_2 \pm P_3)}{\sqrt{2} \Gamma_b(Q)^3 \prod_{j=1}^{3} \Gamma_b(Q\pm 2P_j)}\,.
    \end{aligned}
\end{equation}

There is an important consistency check. The identity limit of the 3 point structure constant in spacelike Liouville should be consistent with the 2 point structure constant. We have
\begin{equation}
    \underset{P_3 \to Q/2}{\lim} C_0^{(b)}(P_1,P_2,P_3) = \frac{1}{\rho_0^{(b)}(P_1)}\left(\delta(P_1-P_2)+\delta(P_1+P_2)\right)\,.
\end{equation}

\paragraph{Timelike case}
We choose a normalzation for which timelike Liouville CFT data is barely the inverse of the spacelike one. We have
\begin{equation}
    \begin{aligned}
        \hat{B}(P_1) &:= \frac{1}{\hat{\rho}_0(P_1)} = \frac{2 \rho_0(iP_1)}{P^2_1}\,,\\
        \hat{C}_0(P_1,P_2,P_3) &= \frac{1}{C_0^{(b)}(iP_1,iP_2,iP_3)}\,.
    \end{aligned}
\end{equation}


\subsection{Spacelike expressions}
\paragraph{Fusion kernel}
We start by the different parameters defining the quantum modular fusion polynomial of the fusion kernel. We first define $c_{mn}(x)=\cos(2\pi s x)$ and $s_{mn}(x)=\sin(2\pi s x)$ as well as 
\begin{equation}
    a_i=P_i-\frac{m+n}{2s}\,.
\end{equation}
The quantum modular fusion determinant is easily expressed in terms of the determinant of the following Gram matrix. We define the symmetric gram matrix
\begin{equation}
    \mathcal{G}^{(f)}_{m,n}=
    \begin{bmatrix}
    1 & -c_{mn}(a_2) & -c_{mn}(a_3) & c_{mn}(a_s) \\
    -c_{mn}(a_2) & 1 & c_{mn}(a_t) & -c_{mn}(a_1) \\
    -c_{mn}(a_3) & c_{mn}(a_t) & 1 & -c_{mn}(a_4) \\
    c_{mn}(a_s) & -c_{mn}(a_1) & -c_{mn}(a_4) & 1
    \end{bmatrix}
\end{equation}
The quantum modular fusion determinant takes the simple form 
\begin{equation}
    D^{(f)}=\frac{-4\det \mathcal{G}^{(f)}_{m,n}}{s_{mn}(P_s)^2 s_{mn}(P_t)^2}\,.
\end{equation}
This would be interesting to display in the 2D plot the lines where the determinant $D^{(f)}$ is -1 with external momenta fixed and $P_s, P_t$ varying. Finally, we also need the expression of the quantum modular fusion polynomial to get the roots $z_{\pm}^{(f)}$. We have
\begin{equation}
    P_{m,n}^{(f)}(z,\vec{P})=\alpha_{m,n}^{(f)}z^2+\beta_{m,n}^{(f)}z+\gamma_{m,n}^{(f)}\,,
\end{equation}
with
\begin{equation}\label{eq:alpha}
    \begin{aligned}
        \alpha_{m,n}^{(f)}=&+\Bigg[\exp\left(2\pi i s^2\left(\frac{m^{-1}+n^{-1}}{4}+\frac{1}{2s}(P_s+P_t-P_1-P_2-P_3-P_4)\right)\right)\\
        &+\exp\left(2\pi i s^2\left(\frac{m^{-1}+n^{-1}}{4}+\frac{1}{2s}(P_s+P_t+P_1+P_2+P_3+P_4)\right)\right)\\
        &+\exp\left(2\pi i s^2\left(\frac{m^{-1}+n^{-1}}{4}+\frac{1}{2s}(-P_s-P_t-P_1+P_2-P_3+P_4)\right)\right)\\
        &+\exp\left(2\pi i s^2\left(\frac{m^{-1}+n^{-1}}{4}+\frac{1}{2s}(-P_s-P_t+P_1-P_2+P_3-P_4)\right)\right) \Bigg]\\
        &-\Bigg[\exp\left(2\pi i s^2\left(-\frac{m^{-1}+n^{-1}}{4}+\frac{1}{2s}(-P_s+P_t-P_1-P_2+P_3+P_4)\right)\right)\\
        &+\exp\left(2\pi i s^2\left(-\frac{m^{-1}+n^{-1}}{4}+\frac{1}{2s}(-P_s+P_t+P_1+P_2-P_3-P_4)\right)\right)\\
        &+\exp\left(2\pi i s^2\left(-\frac{m^{-1}+n^{-1}}{4}+\frac{1}{2s}(P_s-P_t-P_1+P_2+P_3-P_4)\right)\right)\\
        &+\exp\left(2\pi i s^2\left(-\frac{m^{-1}+n^{-1}}{4}+\frac{1}{2s}(P_s-P_t+P_1-P_2-P_3+P_4)\right)\right) \Bigg]\,,\\
    \end{aligned}
\end{equation}
and 
\begin{equation}
    \beta_{m,n}^{(f)}=4\left[s_{mn}(P_1)s_{mn}(P_3)+s_{mn}(P_2)s_{mn}(P_4)+s_{mn}(P_s)s_{mn}(P_t)\right]\,.
\end{equation}
The expression of $\gamma_{m,n}^{(f)}$ is  the complex conjugate of $\alpha_{m,n}^{(f)}$ when the momenta are real. 
\begin{equation}
    \begin{aligned}
        \gamma_{m,n}^{(f)}=&+\Bigg[\exp\left(2\pi i s^2\left(\frac{m^{-1}+n^{-1}}{4}+\frac{1}{2s}(P_s+P_t-P_1-P_2-P_3-P_4)\right)\right)\\
        &+\exp\left(-2\pi i s^2\left(\frac{m^{-1}+n^{-1}}{4}+\frac{1}{2s}(P_s+P_t+P_1+P_2+P_3+P_4)\right)\right)\\
        &+\exp\left(-2\pi i s^2\left(\frac{m^{-1}+n^{-1}}{4}+\frac{1}{2s}(-P_s-P_t-P_1+P_2-P_3+P_4)\right)\right)\\
        &+\exp\left(-2\pi i s^2\left(\frac{m^{-1}+n^{-1}}{4}+\frac{1}{2s}(-P_s-P_t+P_1-P_2+P_3-P_4)\right)\right) \Bigg]\\
        &-\Bigg[\exp\left(-2\pi i s^2\left(-\frac{m^{-1}+n^{-1}}{4}+\frac{1}{2s}(-P_s+P_t-P_1-P_2+P_3+P_4)\right)\right)\\
        &+\exp\left(-2\pi i s^2\left(-\frac{m^{-1}+n^{-1}}{4}+\frac{1}{2s}(-P_s+P_t+P_1+P_2-P_3-P_4)\right)\right)\\
        &+\exp\left(-2\pi i s^2\left(-\frac{m^{-1}+n^{-1}}{4}+\frac{1}{2s}(P_s-P_t-P_1+P_2+P_3-P_4)\right)\right)\\
        &+\exp\left(-2\pi i s^2\left(-\frac{m^{-1}+n^{-1}}{4}+\frac{1}{2s}(P_s-P_t+P_1-P_2-P_3+P_4)\right)\right) \Bigg]\,,\\
    \end{aligned}
\end{equation}
We will not use it to compute the roots of the polynomials but it might serve as a consistency check. Indeed, we have
\begin{equation}
    z_{\pm}^{(f)}=\frac{-\beta_{m,n}^{(f)}\pm2i s_{mn}(P_t)s_{mn}(P_s)\sqrt{D^{(f)}}}{2\alpha_{m,n}^{(f)}}=\frac{-\beta_{m,n}^{(f)}\pm\sqrt{\Delta^{(f)}}}{2\alpha_{m,n}^{(f)}}\,.
\end{equation}
We barely have everything.

We define the $\pm$ kernels whose expressions for rational charge are given in your paper Ioannis https://arxiv.org/pdf/2512.03172.
\begin{equation}\label{eq:elementaryblocks}
    \mathbb{F}^{(\pm)}_{P_s, P_t} \begin{bmatrix}
     P_2 & P_3 \\
     P_1 & P_4
    \end{bmatrix}
    =\pm \frac{F_b(P_1,P_2,P_3,P_4|P_s,P_t)}{2 i \sin(2\pi s P_t)\sin(2\pi s P_s)}\frac{1}{\sqrt{D^{(f)}}}\sum_{k=0}^{s^2-1} \mathcal{F}_b\left(\frac{\log z_{\pm}^{(f)}}{2\pi i s^2}+\frac{k}{s^2}\right)\,.
\end{equation}
We have
\begin{equation}
    \mathcal{F}_b(u)=\frac{\rm Numerator(\mathcal{F}_b)}{\rm Denominator(\mathcal{F}_b)}\,,
\end{equation}
with 
\begin{equation}
\begin{aligned}
\rm Numerator(\mathcal{F}_b)&=\tilde{G}_{m,n}\left(u-1+\frac{m^{-1}+n^{-1}}{4}-\frac{1}{2s}\left(-P_s+P_t-P_1-P_2+P_3+P_4\right)\right)\\
&\times \tilde{G}_{m,n}\left(u-1+\frac{m^{-1}+n^{-1}}{4}-\frac{1}{2s}\left(-P_s+P_t+P_1+P_2-P_3-P_4\right)\right)\\
&\times \tilde{G}_{m,n}\left(u-1+\frac{m^{-1}+n^{-1}}{4}-\frac{1}{2s}\left(P_s-P_t-P_1+P_2+P_3-P_4\right)\right)\\
&\times \tilde{G}_{m,n}\left(u-1+\frac{m^{-1}+n^{-1}}{4}-\frac{1}{2s}\left(P_s-P_t+P_1-P_2-P_3+P_4\right)\right)\,,\\
\end{aligned}
\end{equation}
and 
\begin{equation}
\begin{aligned}
\rm Denominator(\mathcal{F}_b)&=\tilde{G}_{m,n}\left(u-\frac{m^{-1}+n^{-1}}{4}-\frac{1}{2s}\left(P_s+P_t-P_1-P_2-P_3-P_4\right)\right)\\
&\times \tilde{G}_{m,n}\left(u-\frac{m^{-1}+n^{-1}}{4}-\frac{1}{2s}\left(P_s+P_t+P_1+P_2+P_3+P_4\right)\right)\\
&\times \tilde{G}_{m,n}\left(u-\frac{m^{-1}+n^{-1}}{4}-\frac{1}{2s}\left(-P_s-P_t-P_1+P_2-P_3+P_4\right)\right)\\
&\times \tilde{G}_{m,n}\left(u-\frac{m^{-1}+n^{-1}}{4}-\frac{1}{2s}\left(-P_s-P_t+P_1-P_2+P_3-P_4\right)\right)\,.
\end{aligned}
\end{equation}
Next, we also have 
\begin{equation}
    F_b(u)=\frac{\rm Numerator(F_b)}{\rm Denominator(F_b)}\,,
\end{equation}
with
\begin{equation}
\begin{aligned}
    \rm Numerator(F_b)&=i(2\pi)^{4s^2-(m+n)}G_{m,n}\left(\frac{2P_t}{s}\right)G_{m,n}\left(\frac{-2P_t}{s}\right)\\
    &\times G_{m,n}\left(\frac{m^{-1}+n^{-1}}{2}+\frac{1}{s}\left(P_1-P_2+P_s\right)\right)\\
    &\times G_{m,n}\left(\frac{m^{-1}+n^{-1}}{2}+\frac{1}{s}\left(-P_1+P_2+P_s\right)\right)\\
    &\times G_{m,n}\left(\frac{m^{-1}+n^{-1}}{2}+\frac{1}{s}\left(P_1+P_2-P_s\right)\right)\\
    &\times G_{m,n}\left(\frac{m^{-1}+n^{-1}}{2}+\frac{1}{s}\left(-P_1-P_2-P_s\right)\right)\\
    &\times G_{m,n}\left(\frac{m^{-1}+n^{-1}}{2}+\frac{1}{s}\left(P_3-P_4+P_s\right)\right)\\
    &\times G_{m,n}\left(\frac{m^{-1}+n^{-1}}{2}+\frac{1}{s}\left(-P_3+P_4+P_s\right)\right)\\
    &\times G_{m,n}\left(\frac{m^{-1}+n^{-1}}{2}+\frac{1}{s}\left(P_3+P_4-P_s\right)\right)\\
    &\times G_{m,n}\left(\frac{m^{-1}+n^{-1}}{2}+\frac{1}{s}\left(-P_3-P_4-P_s\right)\right)\,,\\
\end{aligned}
\end{equation}
and
\begin{equation}
\begin{aligned}
    \rm Denominator(F_b)&=s\,G_{m,n}\left(m^{-1}+n^{-1}+\frac{2P_s}{s}\right)G_{m,n}\left(m^{-1}+n^{-1}-\frac{2P_s}{s}\right)\\
    &\times G_{m,n}\left(\frac{m^{-1}+n^{-1}}{2}+\frac{1}{s}\left(P_2+P_3+P_t\right)\right)\\
    &\times G_{m,n}\left(\frac{m^{-1}+n^{-1}}{2}+\frac{1}{s}\left(P_2-P_3-P_t\right)\right)\\
    &\times G_{m,n}\left(\frac{m^{-1}+n^{-1}}{2}+\frac{1}{s}\left(-P_2-P_3+P_t\right)\right)\\
    &\times G_{m,n}\left(\frac{m^{-1}+n^{-1}}{2}+\frac{1}{s}\left(-P_2+P_3-P_t\right)\right)\\
    &\times G_{m,n}\left(\frac{m^{-1}+n^{-1}}{2}+\frac{1}{s}\left(P_1+P_4+P_t\right)\right)\\
    &\times G_{m,n}\left(\frac{m^{-1}+n^{-1}}{2}+\frac{1}{s}\left(P_1-P_4-P_t\right)\right)\\
    &\times G_{m,n}\left(\frac{m^{-1}+n^{-1}}{2}+\frac{1}{s}\left(-P_1-P_4+P_t\right)\right)\\
    &\times G_{m,n}\left(\frac{m^{-1}+n^{-1}}{2}+\frac{1}{s}\left(-P_1+P_4-P_t\right)\right)\,.\\
\end{aligned}
\end{equation}

\paragraph{Consistency chekcs for the fusion kernel}

One non trivial check that we can try to verify is the following symmetry of $\mathcal{F}_{b}(u)$ namely that is invariant under $P_s\rightarrow P_3,P_3\rightarrow P_s,P_t\rightarrow P_1,P_1\rightarrow P_t$. A second non trivial is the fact that the sum defined in the $\pm$ kernels is single-valued and thus does not depend on the branch of the logarithm i.e
\begin{equation}
    \sum_{k=0}^{s^2-1} \mathcal{F}_b\left(\frac{\log z_{\pm}^{(f)}}{2\pi i s^2}+\frac{k}{s^2}\right)=\sum_{k=0}^{s^2-1} \mathcal{F}_b\left(\frac{\log z_{\pm}^{(f)}}{2\pi i s^2}+\frac{k+\ell}{s^2}\right)\,,
\end{equation}
for all $\ell$ in $\mathbb{Z}$.

A simple but non trivial property of the kernels is their transformation rules under reflection of internal momenta. Basically, $\mathbb{F}^{(\pm)}$ should get exchanged under reflection of the internal momenta eg $P_s\rightarrow-P_s$ in the spacelike case.

A highly non trivial property is the mixing property which we write
\begin{equation}\label{eq:mixingproperty}
    \mathbb{F}_{P_s,P_t}^{(\eta)}\begin{bmatrix}
     P_2 & P_3 \\
     P_1 & P_4
    \end{bmatrix}=\frac{\mathbb{F}_{\mathbb{I},P_t}^{(TV)}\begin{bmatrix}
     P_2 & P_3 \\
     P_2 & P_3
    \end{bmatrix}}{\mathbb{F}_{\mathbb{I},P_1}^{(TV)}\begin{bmatrix}
     P_2 & P_s \\
     P_2 & P_s
    \end{bmatrix}}\,\mathbb{F}_{P_3,P_1}^{(\eta)}\begin{bmatrix}
     P_2 & P_s \\
     P_t & P_4
    \end{bmatrix}\,,
\end{equation}
where TV stands for the Teschner/Vartanov kernel.
Before checking the pentagon equation for the spacelike kernels, we will first verify a relation that was written by Ioannis and Sylvain in the paper concerning the $c=25$ fusion kernels. Precisely, the relation reads
\begin{equation}\label{eq:RibaultTsiares}
    \mathbb{F}^{(\pm)}_{P_s,P_t}\begin{bmatrix}
     1/4 & 1/4 \\
     1/4 & 1/4
    \end{bmatrix}=\pm i \frac{P_t}{P_s}16^{P_t^2-P_s^2}e^{\pm 2\pi i P_s P_t}\,.
\end{equation}
After checking that we are ready to implement the pentagon identity for the spacelike kernels. In order to test if it satisfied also for $F^{(\pm)}$ and only for the Teschner-Vartanov Kernel, we will implement the following:
\begin{equation}\label{eq:Pentagoninlittlepgenericrelation}
    \begin{aligned}
        &\int_{\mathbb{R}} dp_s \, \mathbb{F}_{P_q, P_4} 
    \begin{bmatrix}
     ip_s & P_3 \\
     P_u & P_5
     \end{bmatrix}
\mathbb{F}_{P_o, ip_s} 
    \begin{bmatrix}
     P_1 & P_2 \\
     P_u & P_q
     \end{bmatrix}
\mathbb{F}^{(TV)}_{ip_s, P_t} 
    \begin{bmatrix}
     P_2 & P_3 \\
     P_1 & P_4
     \end{bmatrix}\\
&= \mathbb{F}_{P_o, P_4} 
    \begin{bmatrix}
     P_1 & P_t \\
     P_u & P_5
     \end{bmatrix}
\mathbb{F}_{P_q, P_t} 
    \begin{bmatrix}
     P_2 & P_3 \\
     P_o & P_5
     \end{bmatrix}
    \end{aligned}\,,
\end{equation}
where $\mathbb{F}$ denotes a generic kernel (it could be $\mathbb{F}^{(\pm)}$ or the Teschner-Vartanov one $\mathbb{F}^{(TV)}$ or something else that has the same arguments).
Finally, the idemptency relation reads
\begin{equation}\label{eq:idempotencyf}
    \int_{\mathbb{R}} dp \,\mathbb{F}^{(b)}_{P_s, ip} \begin{bmatrix}
     P_2 & P_3 \\
     P_1 & P_4
    \end{bmatrix}
    \mathbb{F}^{(b)}_{iP, P_s'} \begin{bmatrix}
     P_3 & P_4 \\
     P_2 & P_1
    \end{bmatrix}
    = \frac{1}{2} \left(\delta(P_s - P_s')+\delta(P_s+P_s')\right)\,.
\end{equation}

\paragraph{Modular kernel}
We stary by defining the necessary expressions for the quantum modular modular polynomial. Since you give explicit expressions in the article, I will directly use these expressions. The quantum modular modular determinant reads
\begin{equation}
    \mathcal{D}^{(m)}=1-\left(\frac{s_{mn}\left(\frac{P_0}{2}-\frac{m+n}{4s}\right)}{s_{mn}(P_t)s_{mn}(P_s)}\right)^2\,,
\end{equation}
while the roots are 
\begin{equation}
    z_{\pm}^{(m)}=\frac{s_{mn}(P_s)}{s_{mn}\left(\frac{P_0}{2}-P_s-\frac{m+n}{4s}\right)}\left[(-1)^{mn}c_{mn}(P_t)\mp i s_{mn}(P_t)\sqrt{\mathcal{D}^{(m)}}\right]\,.
\end{equation}
We define the $\pm$ kernels whose expressions for rational charge are given in your paper Ioannis https://arxiv.org/pdf/2512.03172.
\begin{equation}\label{eq:elementaryblocks}
    \mathbb{M}^{(\pm)}_{P_s, P_t} [P_0]
    =\pm \frac{M_b(P_0|P_s,P_t)}{2 i \sin(2\pi s P_t)\sin(2\pi s P_s)}\frac{1}{\sqrt{D^{(m)}}}\sum_{k=0}^{s^2-1} \mathcal{M}_b\left(\frac{\log z_{\pm}^{(m)}}{2\pi i s^2}+\frac{k}{s^2}\right)\,.
\end{equation}
We have
\begin{equation}
    \mathcal{M}_{b}(u)=\frac{\rm {Numerator} (\mathcal{M}_{b})}{\rm {Denominator} (\mathcal{M}_{b})}
\end{equation}
with
\begin{equation}
    \begin{aligned}
        \rm {Numerator} (\mathcal{M}_{b})&=e^{-4\pi i s P_s u}\tilde{G}_{m,n}\left(u-1+\frac{m^{-1}+n^{-1}}{4}-\frac{1}{2}\left(\frac{2P_t}{s}+\frac{P_0}{s}\right)\right)\\
        &\times\tilde{G}_{m,n}\left(u-1+\frac{m^{-1}+n^{-1}}{4}-\frac{1}{2}\left(-\frac{2P_t}{s}+\frac{P_0}{s}\right)\right)
    \end{aligned}
\end{equation}
and
\begin{equation}
    \begin{aligned}
        \rm {Denominator} (\mathcal{M}_{b})&=\tilde{G}_{m,n}\left(u-\frac{m^{-1}+n^{-1}}{4}-\frac{1}{2}\left(\frac{2P_t}{s}-\frac{P_0}{s}\right)\right)\\
        &\times\tilde{G}_{m,n}\left(u-\frac{m^{-1}+n^{-1}}{4}-\frac{1}{2}\left(-\frac{2P_t}{s}-\frac{P_0}{s}\right)\right)
    \end{aligned}
\end{equation}
Next, we also have 
\begin{equation}
    M_b(P_0|P_s,P_t)=\frac{\rm Numerator(M_b)}{\rm Denominator(M_b)}\,,
\end{equation}
with 
\begin{equation}
    \begin{aligned}
        \rm Numerator(M_b)&=s (2\pi)^{sP_0-m-n}\mathcal{N}_{m,n}(P_s)\rho_0^{(b)}(P_t)\\
        &\times G_{m,n}\left(m^{-1}+n^{-1}+\frac{2P_t}{s}\right)G_{m,n}\left(\frac{m^{-1}+n^{-1}}{2}-\frac{P_0+2P_s}{s}\right)\\
        &\times G_{m,n}\left(m^{-1}+n^{-1}-\frac{2P_t}{s}\right)G_{m,n}\left(\frac{m^{-1}+n^{-1}}{2}-\frac{P_0-2P_s}{s}\right)\,,
    \end{aligned}
\end{equation}
with $\mathcal{N}_{n,m}=(-1)^{nm}(2\pi)^{nm}s^{-2}e^{2\pi i s P_s}$ and $\rho_0^{(b)}$ is the Plancherel measure 
\[\rho_0^{(b)}(P_t)=-4\sqrt{2}\sin(2\pi b P_t)\sin(2\pi b^{-1} P_t)\,.\]
Finally the denominator reads
\begin{equation}
    \begin{aligned}
        \rm Denominator(M_b)&=2\tilde{G}_{m,n}\left(\frac{-P_0}{s}\right)\\
        &\times G_{m,n}\left(m^{-1}+n^{-1}+\frac{2P_s}{s}\right)G_{m,n}\left(\frac{m^{-1}+n^{-1}}{2}-\frac{P_0+2P_t}{s}\right)\\
        &\times G_{m,n}\left(m^{-1}+n^{-1}-\frac{2P_s}{s}\right)G_{m,n}\left(\frac{m^{-1}+n^{-1}}{2}-\frac{P_0-2P_t}{s}\right)\,.
    \end{aligned}
\end{equation}

\paragraph{Consistency checks for the modular kernel} Before checking the idempotency relation and the torus two point function consistency condition, we will start by some trivial checks which will make sure that we implemented the correct definition. First, the sum defined in the $\pm$ kernels is single-valued and thus does not depend on the branch of the logarithm i.e
\begin{equation}
    \sum_{k=0}^{s^2-1} \mathcal{M}_b\left(\frac{\log z_{\pm}^{(f)}}{2\pi i s^2}+\frac{k}{s^2}\right)=\sum_{k=0}^{s^2-1} \mathcal{M}_b\left(\frac{\log z_{\pm}^{(f)}}{2\pi i s^2}+\frac{k+\ell}{s^2}\right)\,,
\end{equation}
for all $\ell$ in $\mathbb{Z}$. A second non trivial check is the identity limit of the Teschner-Vartanov modular kernel precisely,
\begin{equation}\label{eq:identitylimit}
    \mathbb{M}^{(TV)}_{P_s,P_t}\left[P_0=-\frac{b+b^{-1}}{2}\right]=\sqrt{2}\cos(4\pi P_s P_t)\,.
\end{equation}
It captures modular covariance of the vacuum representation. Finally, we would like to check the reflection properties of the $\mathbb{M}^{(\pm)}_{P_s,P_t}\left[P_0\right]$ kernels. They have to transform into one another under reflection of both internal momenta, precisely,
\begin{equation}\label{eq:reflectionproperty}
    \begin{aligned}
        \mathbb{M}^{(\pm)}_{P_s,P_t}\left[P_0\right] &\underset{P_s\to -P_s}{\rightarrow} \mathbb{M}^{(\mp)}_{P_s,P_t}\left[P_0\right]\\
        \mathbb{M}^{(\pm)}_{P_s,P_t}\left[P_0\right] &\underset{P_t\to -P_t}{\rightarrow} \mathbb{M}^{(\mp)}_{P_s,P_t}\left[P_0\right]\,.
    \end{aligned}
\end{equation}
Instead of checking the idempotency relation for the $\pm$ kernels, we will check the following relation which is a consequence of the idempotency relation
\begin{equation}
    \mathbb{M}_{P_1,P'}[P_0] = \frac{\frac{C_0(P', P', P_0)}{B(P')}}{\frac{C_0(P_1, P_1, P_0)}{B(P_1)}}\,\mathbb{M}_{P', P_1}[P_0]\,.
\end{equation}
We did not specify if this concerns the $\pm$ kernels or the TV kernel because we would like to check it for both of them.

\paragraph{Consistency checks involving the Fusion and the modular kernels}
Before turning to the most difficult equation which is consistency of the torus two point functions, we have the non-rational Verlinde formula presented in your paper Ioannis.
\begin{equation}
    \begin{aligned}
        \mathbb{F}^{(b)}_{P_1,P_3}\begin{bmatrix}
            P_0 & P_3 \\
            P_1 & P_2
        \end{bmatrix}
        = \int_{\mathbb{R}} dp \frac{\mathbb{M}_{P_1,ip}[P_0]\,\mathbb{M}_{P_2,ip}[\mathbb{I}]\,\mathbb{M}_{ip, P_3}[P_0]}{\mathbb{M}_{\mathbb{I},ip}[\mathbb{I}]}\,.
    \end{aligned}
\end{equation}
Finally, we have the consistency condition for the torus two point function which reads
\begin{equation}
    \begin{aligned}
        \mathbb{M}_{P_1,P_2}[P_3]&\int_{\mathbb{R}} dp_4 \, \mathbb{F}_{P_3, ip_4}\begin{bmatrix}
            P_0' & P_2 \\
            P_0 & P_2
        \end{bmatrix}
        e^{\pi i (2h_4-2h_2+\frac{h_3}{2})}\mathbb{F}_{ip_4, P_5} \begin{bmatrix}
            P_0' & P_0 \\
            P_2 & P_2
        \end{bmatrix}\\
        =& \int_{\mathbb{R}} dp_6 \, \mathbb{F}_{P_3, ip_6}\begin{bmatrix}
            P_0 & P_1 \\
            P_0' & P_1
        \end{bmatrix}
        \mathbb{F}_{P_1,P_5}\begin{bmatrix}
            P_0' & P_0 \\
            ip_6 & ip_6
        \end{bmatrix}
        e^{\pi i (h_0+h_0'-\frac{h_5}{2})}\mathbb{M}_{ip_6,P_2}[P_5]\,.
    \end{aligned}
\end{equation}
As the LHS integral is not absolutely convergent, we wil shift the contour as
\begin{equation}
    \begin{aligned}
        \mathbb{M}_{P_1,P_2}[P_3]&\int_{\mathbb{R}} d(p_4e^{i\pi /12}) \, \mathbb{F}_{P_3, i p_4 e^{i\pi /12}}\begin{bmatrix}
            P_0' & P_2 \\
            P_0 & P_2
        \end{bmatrix}
        e^{\pi i (2h_4(p_4e^{i\pi /12})-2h_2+\frac{h_3}{2})}\mathbb{F}_{ip_4 e^{i\pi /12}, P_5} \begin{bmatrix}
            P_0' & P_0 \\
            P_2 & P_2
        \end{bmatrix}\\
        =& \int_{\mathbb{R}} dp_6 \, \mathbb{F}_{P_3, ip_6}\begin{bmatrix}
            P_0 & P_1 \\
            P_0' & P_1
        \end{bmatrix}
        \mathbb{F}_{P_1,P_5}\begin{bmatrix}
            P_0' & P_0 \\
            ip_6 & ip_6
        \end{bmatrix}
        e^{\pi i (h_0+h_0'-\frac{h_5}{2})}\mathbb{M}_{ip_6,P_2}[P_5]\,.
    \end{aligned}
\end{equation}

\subsection{Timelike expressions}

\paragraph{Virasoro Wick rotation}
The kernels are obtained through the Virasoro Wick rotation of the spacelike kernels. It reads
\begin{equation}\label{eq:VWR}
    \begin{aligned}
        \hat{\mathbb{M}}_{P_s,P_t}^{(\pm)}[P_0] &= \mp i \frac{P_t}{P_s}\mathbb{M}_{iP_t,iP_s}[iP_0] \\
        \hat{\mathbb{F}}_{P_s,P_t}^{(\pm)}\begin{bmatrix}
     P_2 & P_3 \\
     P_1 & P_4
    \end{bmatrix} &= \mp i \frac{P_t}{P_s} \mathbb{F}_{iP_t,iP_s}^{(\pm)}\begin{bmatrix}
     iP_2 & iP_1 \\
     iP_3 & iP_4
    \end{bmatrix}
    \end{aligned}
\end{equation}
Note that the value of $b$ remains formally unchanged that is to obtain the timelike kernel at $c=1$, it suffices to apply the prescription~\eqref{eq:VWR} with $b=1$ for evaluation of the RHS.

\paragraph{Consistency chekcs for the fusion kernel}

A simple but non trivial property of the kernels is their transformation rules under reflection of internal momenta. Basically, $\hat{\mathbb{F}}^{(\pm)}$ should get exchanged under reflection of the internal momenta eg $P_s\rightarrow-P_s$. Before checking the pentagon equation for the spacelike kernels, we will first verify a relation that was written by Ioannis and Sylvain in the paper concerning the $c=1$ fusion kernels. Precisely, the relation reads
\begin{equation}\label{eq:RibaultTsiaresc=1}
    \hat{\mathbb{F}}^{(\pm)}_{P_s,P_t}\begin{bmatrix}
     i/4 & i/4 \\
     i/4 & i/4
    \end{bmatrix}=16^{P_t^2-P_s^2}e^{\mp 2\pi i P_s P_t}\,.
\end{equation}
We also have the analog of the mixing symmetry for the timelike fusion kernels which reads
\begin{equation}
    \hat{\mathbb{F}}_{P_s,P_t}^{(\eta)}\begin{bmatrix}
     P_2 & P_3 \\
     P_1 & P_4
    \end{bmatrix}=\frac{\Psi_{P_3, P_t}[P_2]}{\Psi_{P_s, P_1}[P_2]}\,\hat{\mathbb{F}}_{P_3,P_1}^{(\eta)}\begin{bmatrix}
     P_2 & P_s \\
     P_t & P_4
    \end{bmatrix}\,,
\end{equation}
where we define 
\begin{equation}
    \Psi_{P_1, P_2}[P_3]=\frac{P_2}{P_1}\hat{\rho}_0(P_1)\hat{C}_0(P_1, P_2, P_3)\,.
\end{equation}
After checking that we are ready to implement the pentagon identity for the spacelike kernels. In order to test if it satisfied also for $\hat{\mathbb{F}}^{(\pm)}$ and only for the analog of the TV kernel, we will implement the following:
\begin{equation}\label{eq:Pentagoninlittlepgenericrelation}
    \begin{aligned}
        &\int_{\mathbb{R}} dp_s \, \hat{\mathbb{F}}_{P_q, P_4} 
    \begin{bmatrix}
     ip_s & P_3 \\
     P_u & P_5
     \end{bmatrix}
\hat{\mathbb{F}}_{P_o, ip_s} 
    \begin{bmatrix}
     P_1 & P_2 \\
     P_u & P_q
     \end{bmatrix}
\hat{\mathbb{F}}^{(\rm {phys.})}_{ip_s, P_t} 
    \begin{bmatrix}
     P_2 & P_3 \\
     P_1 & P_4
     \end{bmatrix}\\
&= \hat{\mathbb{F}}_{P_o, P_4} 
    \begin{bmatrix}
     P_1 & P_t \\
     P_u & P_5
     \end{bmatrix}
\hat{\mathbb{F}}_{P_q, P_t} 
    \begin{bmatrix}
     P_2 & P_3 \\
     P_o & P_5
     \end{bmatrix}
    \end{aligned}\,,
\end{equation}
where $\hat{\mathbb{F}}$ denotes a generic kernel (it could be $\hat{\mathbb{F}^{(\pm)}}$ or the physical one $\hat{\mathbb{F}}^{(\rm {phys.})}$ or something else that has the same arguments).
Finally, the idemptency relation reads
\begin{equation}\label{eq:idempotencyf}
    \int_{\mathbb{R}} dp \,\hat{\mathbb{F}}^{(b)}_{P_s, ip} \begin{bmatrix}
     P_2 & P_3 \\
     P_1 & P_4
    \end{bmatrix}
    \hat{\mathbb{F}}^{(b)}_{iP, P_s'} \begin{bmatrix}
     P_3 & P_4 \\
     P_2 & P_1
    \end{bmatrix}
    = \frac{1}{2} \left(\delta(P_s - P_s')+\delta(P_s+P_s')\right)\,.
\end{equation}


\section{Numerical Packages/conventions}
Since the only library that I found that does implement Barnes G-function is mpmath, we will use mpmath for the special functions.
\begin{itemize}[leftmargin=*]
    \item Error estimates / tolerances (SciPy quad, etc.)
    \item Complex integration conventions (parametrized contours, branch choices)
\end{itemize}

\section{Explanations/Workflow}
I invite you also to check the end of the workflow since there are a bunch of checks that everythings work correctly so far.
\paragraph{Pentagon equation/spacelike}
The values of $P$'s and $m,n$ will be our global variables as well as the cutoff $\Lambda$ for the integration. We start by the packages and we define a dataclass for our global variables. Also, we define a DEMO mod which is True for the demo tests and false for the final checks which we will perform at the very end.
\begin{lstlisting}[language=Python]
#Packages
from dataclasses import dataclass
from mpmath import mp, barnesg
import numpy as np
import matplotlib.pyplot as plt

# -------- Global settings --------
mp.dps = 20 #be careful with the precision, if the precision is not enough you might raise "Zero Division error"
DEMO = True  # True for Demo mode

# -------- Parameters container --------
@dataclass(frozen=True)
class ParametersPentagon:
    Po: mp.mpc
    Pq: mp.mpc
    P1: mp.mpc
    P2: mp.mpc
    P3: mp.mpc
    P4: mp.mpc
    P5: mp.mpc
    Pt: mp.mpc
    Pu: mp.mpc
@dataclass(frozen=True)
class ParametersIdempotency:
    Ps: mp.mpc
    Psprime: mp.mpc
    P1: mp.mpc
    P2: mp.mpc
    P3: mp.mpc
    P4: mp.mpc
    
    
@dataclass
class Theory:
    m: int
    n: int
    Lambda: mp.mpf = mp.mpf(3) # cutoff (for later integration)

    @property
    def b(self):
        return mp.sqrt(mp.mpf(self.m) / mp.mpf(self.n))

    @property
    def s(self):
        return mp.sqrt(mp.mpf(self.m) * mp.mpf(self.n))


# -------- Demo / main params --------
if DEMO:
    parPenta = ParametersPentagon(
        Po=mp.mpc(0.001+0.1j), Pq=mp.mpc(0.001+0.5j), 
        P1=mp.mpc(0.001+0.2j), P2=mp.mpc(0.001+0.6j), 
        P3=mp.mpc(0.001+0.3j), P4=mp.mpc(0.001+0.7j), 
        P5=mp.mpc(0.001+0.4j), Pt=mp.mpc(0.001+0.8j),
        Pu=mp.mpc(0.001+0.9j)
    )
    parIdem = ParametersIdempotency(
        Ps=mp.mpc(0.001+0.1j), Psprime=mp.mpc(0.001+3j), 
        P1=mp.mpc(0.001+0.2j), P2=mp.mpc(0.001+0.6j), 
        P3=mp.mpc(0.001+0.3j), P4=mp.mpc(0.001+0.7j),
    )
    theo=Theory(m=1,n=1)
else:
    parPenta = ParametersPentagon(
        Po=mp.mpc(0.0+0.1j), Pq=mp.mpc(0.0+0.2j),
        P1=mp.mpc(0.0+0.3j), P2=mp.mpc(0.0+0.4j),
        P3=mp.mpc(0.0+0.5j), P4=mp.mpc(0.0+0.6j),
        P5=mp.mpc(0.0+0.7j), Pt=mp.mpc(0.0+0.8j),
        Pu=mp.mpc(0.001+0.9j)
    )
    theo=Theory(m=1,n=1)

\end{lstlisting}
In addition to these preliminary data that we defined, we also introduce some basics modules that could be useful if we want to plot the phase or the modulus of a given function.
\begin{lstlisting}
# =========================================
# Useful modules to plot a bunch of different things
# =========================================

def plot_phase2D(func, x_min=-2, x_max=2, y_min=-2, y_max=2, resolution=100):

   x = np.linspace(x_min, x_max, resolution)
   y = np.linspace(y_min, y_max, resolution)
   X, Y = np.meshgrid(x, y)
   Z = np.zeros_like(X, dtype=np.float64)

   for i in range(resolution):
       for j in range(resolution):
           z_val = func(mp.mpc(X[i, j], Y[i, j]))
           Z[i, j] = mp.arg(z_val)

   plt.figure(figsize=(8, 6))
   plt.contourf(X, Y, Z, levels=100, cmap='hsv')
   plt.colorbar(label='Phase (radians)')
   plt.title('Phase of the function in the complex plane')
   plt.xlabel('Real part')
   plt.ylabel('Imaginary part')
   plt.xlim(x_min, x_max)
   plt.ylim(y_min, y_max)
   plt.show()

def plot_module1D(func, name: str, x_min=-5, x_max=5, resolution=300):
    x = np.linspace(x_min, x_max, resolution)
    y = np.zeros(resolution, dtype=float)

    for i in range(resolution):
        z_val = func(mp.mpc(x[i]))
        y[i] = float(abs(z_val))

    plt.figure(figsize=(8, 6))
    plt.plot(x, y)
    plt.title('Modulus of the function')
    plt.xlabel('p')
    plt.ylabel(name)
    plt.xlim(x_min, x_max)
    plt.grid(True)
    plt.show()

def plot_real1D(func, name: str, x_min=-5, x_max=5, resolution=300):
    x = np.linspace(x_min, x_max, resolution)
    y = np.zeros(resolution, dtype=float)

    for i in range(resolution):
        z_val = func(mp.mpc(x[i]))
        y[i] = float(mp.re(z_val))

    plt.figure(figsize=(8, 6))
    plt.plot(x, y)
    plt.title('Real part of the function')
    plt.xlabel('p')
    plt.ylabel(name)
    plt.xlim(x_min, x_max)
    plt.grid(True)
    plt.show()
\end{lstlisting}
We then define the special functions $\tilde{G}$, $G_{m,n}$ and $\tilde{G}_{m,n}$ (cf~\eqref{eq:Gtilde},~\eqref{eq:Gmn} and~\eqref{eq:Gtildemn}). These functions have some properties that we can check numerically (cf~\eqref{eq:ComplexConjugation} and~\eqref{eq:InversionRelations}). We will also check the shift relation for $G_{m,n}$ (cf~\eqref{eq:CheckGmn}). These checks will be implemented at the end of the workflow.
\begin{lstlisting}[language=Python]
# =========================================
# Barnes G products
# =========================================

def Gmn(Theory: Theory, z:complex):
    m, n = Theory.m, Theory.n
    z = mp.mpc(z)

    prod = mp.mpc(1)
    m_mp = mp.mpf(m)
    n_mp = mp.mpf(n)

    for k in range(m):
        kshift = mp.mpf(k) / m_mp
        for l in range(n):
            prod *= barnesg(z + kshift + mp.mpf(l) / n_mp)
    return prod


def Gtilde(z:complex):
    """
    \tilde G(z) = G(1+z)/G(1-z)
    """
    z = mp.mpc(z)
    return barnesg(1 + z) / barnesg(1 - z)


def Gtilde_mn(Theory: Theory, z:complex):
    m, n = Theory.m, Theory.n
    z = mp.mpc(z)

    prod = mp.mpc(1)
    m_mp = mp.mpf(m)
    n_mp = mp.mpf(n)

    half = mp.mpf(1) / 2
    for k in range(m):
        kshift = mp.mpf(k) / m_mp + half / m_mp
        for l in range(n):
            shift = kshift + mp.mpf(l) / n_mp + half / n_mp
            prod *= Gtilde(z + shift - 1)
    return prod
\end{lstlisting}
We will then turn on to the expression of the quantum modular polynomial and the determinant. First, we implement the trigonometric functions $c_{mn}$ and $s_{mn}$ as well as the shifted momenta $a_i$. 
\begin{lstlisting}[language=Python]
# =========================================
# Trig helpers + a(z)
# =========================================

def smn(Theory: Theory, z:complex):
    return mp.sin(2 * mp.pi * Theory.s * mp.mpc(z))

def cmn(Theory: Theory, z:complex):
    return mp.cos(2 * mp.pi * Theory.s * mp.mpc(z))

def a(Theory: Theory, z:complex):
    # a_i = P_i - (m+n)/(2s) in your notation; here: a(Theory, z) = z - (m+n)/(2s)
    return mp.mpc(z) - (mp.mpf(Theory.m + Theory.n) / (2 * Theory.s))
\end{lstlisting}
We start by the fusion. We first define all the relevant quantities for the quantum modular fusion polynomial
\begin{lstlisting}[language=Python]
# =========================================
# ===================
# FUSION KERNEL
# ===================
# =========================================


# =========================================
# Quantum modular fusion polynomial
# =========================================

# Parameters of the polynomial
def alphaf(Theory: Theory, P1:complex, P2:complex, P3:complex, P4:complex, Ps:complex, Pt:complex):
    
    P1, P2, P3, P4, Ps, Pt = mp.mpc(P1), mp.mpc(P2), mp.mpc(P3), mp.mpc(P4), mp.mpc(Ps), mp.mpc(Pt)
    m, n = Theory.m, Theory.n
    s = Theory.s

    m_inv = mp.mpf(1) / mp.mpf(m)
    n_inv = mp.mpf(1) / mp.mpf(n)

    pref_plus  = 2 * mp.pi * 1j * s**2 * (m_inv + n_inv) / 4
    pref_minus = 2 * mp.pi * 1j * s**2 * (-m_inv - n_inv) / 4

    def E(pref, lin):
        return mp.e ** (pref + mp.pi * 1j * s * lin)

    term1 = E(pref_plus,  Ps + Pt - P1 - P2 - P3 - P4)
    term2 = E(pref_plus,  Ps + Pt + P1 + P2 + P3 + P4)
    term3 = E(pref_plus, -Ps - Pt - P1 + P2 - P3 + P4)
    term4 = E(pref_plus, -Ps - Pt + P1 - P2 + P3 - P4)

    term5 = E(pref_minus, -Ps + Pt - P1 - P2 + P3 + P4)
    term6 = E(pref_minus, -Ps + Pt + P1 + P2 - P3 - P4)
    term7 = E(pref_minus,  Ps - Pt - P1 + P2 + P3 - P4)
    term8 = E(pref_minus,  Ps - Pt + P1 - P2 - P3 + P4)

    return term1 + term2 + term3 + term4 - term5 - term6 - term7 - term8

def betaf(Theory: Theory, P1:complex, P2:complex, P3:complex, P4:complex, Ps:complex, Pt:complex):
    return 4*(smn(Theory,P1)*smn(Theory,P3)+smn(Theory,P2)*smn(Theory,P4)+smn(Theory,Ps)*smn(Theory,Pt))

def gammaf(Theory: Theory, P1:complex, P2:complex, P3:complex, P4:complex, Ps:complex, Pt:complex):
    P1, P2, P3, P4, Ps, Pt = mp.mpc(P1), mp.mpc(P2), mp.mpc(P3), mp.mpc(P4), mp.mpc(Ps), mp.mpc(Pt)
    m, n = Theory.m, Theory.n
    s = Theory.s

    m_inv = mp.mpf(1) / mp.mpf(m)
    n_inv = mp.mpf(1) / mp.mpf(n)

    pref_plus  = 2 * mp.pi * 1j * s**2 * (m_inv + n_inv) / 4
    pref_minus = 2 * mp.pi * 1j * s**2 * (-m_inv - n_inv) / 4

    def E(pref, lin):
        return mp.e ** (-pref - mp.pi * 1j * s * lin)

    term1 = E(pref_plus,  Ps + Pt - P1 - P2 - P3 - P4)
    term2 = E(pref_plus,  Ps + Pt + P1 + P2 + P3 + P4)
    term3 = E(pref_plus, -Ps - Pt - P1 + P2 - P3 + P4)
    term4 = E(pref_plus, -Ps - Pt + P1 - P2 + P3 - P4)

    term5 = E(pref_minus, -Ps + Pt - P1 - P2 + P3 + P4)
    term6 = E(pref_minus, -Ps + Pt + P1 + P2 - P3 - P4)
    term7 = E(pref_minus,  Ps - Pt - P1 + P2 + P3 - P4)
    term8 = E(pref_minus,  Ps - Pt + P1 - P2 - P3 + P4)
    
    return term1 + term2 + term3 + term4 - term5 - term6 - term7 - term8


#Determinant of the quantum modular polynomial
def Gram_matrixf(Theory: Theory, P1:complex, P2:complex, P3:complex, P4:complex, Ps:complex, Pt:complex):
    P1, P2, P3, P4, Ps, Pt = mp.mpc(P1), mp.mpc(P2), mp.mpc(P3), mp.mpc(P4), mp.mpc(Ps), mp.mpc(Pt)
    m, n = Theory.m, Theory.n
    s = Theory.s
    G = mp.zeros(4, 4)
    for i in range(4):
    G[i, i] = 1
    G[0,1]=-cmn(Theory, a(Theory, P2))
    G[0,2]=-cmn(Theory, a(Theory, P3))
    G[0,3]=cmn(Theory, a(Theory, Ps))
    G[1,0]=-cmn(Theory, a(Theory, P2))
    G[1,2]=cmn(Theory, a(Theory, Pt))
    G[1,3]=-cmn(Theory, a(Theory, P1))
    G[2,0]=-cmn(Theory, a(Theory, P3))
    G[2,1]=cmn(Theory, a(Theory, Pt))
    G[2,3]=-cmn(Theory, a(Theory, P4))
    G[3,0]=cmn(Theory, a(Theory, Ps))
    G[3,1]=-cmn(Theory, a(Theory, P1))
    G[3,2]=-cmn(Theory, a(Theory, P4))
    return G

def Gram_determinantf(Theory: Theory, P1:complex, P2:complex, P3:complex, P4:complex, Ps:complex, Pt:complex):
    G = Gram_matrixf(Theory, P1, P2, P3, P4, Ps, Pt)
    return mp.det(G)

def quantum_determinantf(Theory: Theory, P1:complex, P2:complex, P3:complex, P4:complex, Ps:complex, Pt:complex):
    return -4*Gram_determinantf(Theory, P1, P2, P3, P4, Ps, Pt)/(smn(Theory, Ps)**2*smn(Theory, Pt)**2)

# roots of the polynomial
def rootsf(Theory: Theory, P1:complex, P2:complex, P3:complex, P4:complex, Ps:complex, Pt:complex):
    beta_val = betaf(Theory, P1, P2, P3, P4, Ps, Pt)
    alpha_val = alphaf(Theory, P1, P2, P3, P4, Ps, Pt)

    z_plus = (-beta_val +2j*smn(Theory, Ps)*smn(Theory, Pt)*mp.sqrt(quantum_determinantf(Theory, P1, P2, P3, P4, Ps, Pt))) / (2 * alpha_val)
    z_minus = (-beta_val -2j*smn(Theory, Ps)*smn(Theory, Pt)*mp.sqrt(quantum_determinantf(Theory, P1, P2, P3, P4, Ps, Pt))) / (2 * alpha_val)

    return z_plus, z_minus

\end{lstlisting}
We then define the $\pm$ kernels as well as the (TV) kernel. The function "CurlyF" corresponds to $\mathcal{F}_{b}$ while "F" corresponds to $F_{b}$. These are the elementary blocks from which the kernels are built (see~\eqref{eq:elementaryblocks}). The module "Kernels" does implement the kernels. Note that this module contains $\epsilon$ as an argument. $\epsilon=\pm 1$ corresponds to the kernels $F^{\pm}$ while $\epsilon=0$ implements the TV kernel.
\begin{lstlisting}
# =========================================
# Fusion Kernels
# =========================================

# These are the elementary building blocks of the kernels
def CurlyF(Theory: Theory, P1:complex, P2:complex, P3:complex, P4:complex, Ps:complex, Pt:complex, z:complex):
    P1, P2, P3, P4, Ps, Pt, z = mp.mpc(P1), mp.mpc(P2), mp.mpc(P3), mp.mpc(P4), mp.mpc(Ps), mp.mpc(Pt), mp.mpc(z)
    m, n = Theory.m, Theory.n
    s = Theory.s

    term1num = Gtilde_mn(Theory, z -1+(1/m + 1/n)/4 - (-Ps+Pt-P1-P2+P3+P4)/(2*s))
    term2num = Gtilde_mn(Theory, z -1+(1/m + 1/n)/4 - (-Ps+Pt+P1+P2-P3-P4)/(2*s))
    term3num = Gtilde_mn(Theory, z -1+(1/m + 1/n)/4 - (Ps-Pt-P1+P2+P3-P4)/(2*s))
    term4num = Gtilde_mn(Theory, z -1+(1/m + 1/n)/4 - (Ps-Pt+P1-P2-P3+P4)/(2*s))
    Numerator = term1num*term2num*term3num*term4num

    term1den = Gtilde_mn(Theory, z -(1/m + 1/n)/4 - (Ps+Pt-P1-P2-P3-P4)/(2*s))
    term2den = Gtilde_mn(Theory, z -(1/m + 1/n)/4 - (Ps+Pt+P1+P2+P3+P4)/(2*s))
    term3den = Gtilde_mn(Theory, z -(1/m + 1/n)/4 - (-Ps-Pt-P1+P2-P3+P4)/(2*s))
    term4den = Gtilde_mn(Theory, z -(1/m + 1/n)/4 - (-Ps-Pt+P1-P2+P3-P4)/(2*s))
    Denominator = term1den*term2den*term3den*term4den

    return Numerator/Denominator
    
def F(Theory: Theory, P1:complex, P2:complex, P3:complex, P4:complex, Ps:complex, Pt:complex):
    P1, P2, P3, P4, Ps, Pt = mp.mpc(P1), mp.mpc(P2), mp.mpc(P3), mp.mpc(P4), mp.mpc(Ps), mp.mpc(Pt)
    m, n = Theory.m, Theory.n
    s = Theory.s

    term0num=1j*(2*mp.pi)**(4*s**2-(m+n))*Gmn(Theory,2*Pt/s)*Gmn(Theory,-2*Pt/s)
    term1num=Gmn(Theory,(1/m + 1/n)/2 +(P1-P2+Ps)/s)
    term2num=Gmn(Theory,(1/m + 1/n)/2 +(-P1+P2+Ps)/s)
    term3num=Gmn(Theory,(1/m + 1/n)/2 +(P1+P2-Ps)/s)
    term4num=Gmn(Theory,(1/m + 1/n)/2 +(-P1-P2-Ps)/s)
    term5num=Gmn(Theory,(1/m + 1/n)/2 +(P3-P4+Ps)/s)
    term6num=Gmn(Theory,(1/m + 1/n)/2 +(-P3+P4+Ps)/s)
    term7num=Gmn(Theory,(1/m + 1/n)/2 +(P3+P4-Ps)/s)
    term8num=Gmn(Theory,(1/m + 1/n)/2 +(-P3-P4-Ps)/s)
    Numerator = term0num*term1num*term2num*term3num*term4num*term5num*term6num*term7num*term8num

    term0den=s*Gmn(Theory,1/m +1/n + 2*Ps/s)*Gmn(Theory,1/m +1/n - 2*Ps/s)
    term1den=Gmn(Theory,(1/m + 1/n)/2 +(P2+P3+Pt)/s)
    term2den=Gmn(Theory,(1/m + 1/n)/2 +(P2-P3-Pt)/s)
    term3den=Gmn(Theory,(1/m + 1/n)/2 +(-P2-P3+Pt)/s)
    term4den=Gmn(Theory,(1/m + 1/n)/2 +(-P2+P3-Pt)/s)
    term5den=Gmn(Theory,(1/m + 1/n)/2 +(P1+P4+Pt)/s)
    term6den=Gmn(Theory,(1/m + 1/n)/2 +(P1-P4-Pt)/s)
    term7den=Gmn(Theory,(1/m + 1/n)/2 +(-P1-P4+Pt)/s)
    term8den=Gmn(Theory,(1/m + 1/n)/2 +(-P1+P4-Pt)/s)
    Denominator = term0den*term1den*term2den*term3den*term4den*term5den*term6den*term7den*term8den

    return Numerator/Denominator

#These are the plus/minus kernels/TV kernels
def Kernelsf(Theory: Theory, P1: complex, P2: complex, P3: complex, P4: complex, Ps: complex, Pt: complex, epsilon: int):

    # epsilon =\pm are the elementary kernels while epsilon=0 is the Teschner-Vartanov kernel
    s = Theory.s
    N = Theory.m * Theory.n  # this is s^2 as an integer

    zplus, zminus = rootsf(Theory, P1, P2, P3, P4, Ps, Pt)

    prefactor = (
        F(Theory, P1, P2, P3, P4, Ps, Pt)
        / (2j * smn(Theory, Pt) * smn(Theory, Ps) * mp.sqrt(quantum_determinantf(Theory, P1, P2, P3, P4, Ps, Pt)))
    )

    bigsum = mp.mpc(0)

    if epsilon == 1:
        u0 = mp.log(zplus) / (2j * mp.pi * s**2)
        for k in range(N):
            bigsum += CurlyF(Theory, P1, P2, P3, P4, Ps, Pt, u0 + mp.mpf(k) / N)
        return prefactor * bigsum

    elif epsilon == -1:
        u0 = mp.log(zminus) / (2j * mp.pi * s**2)
        for k in range(N):
            bigsum += CurlyF(Theory, P1, P2, P3, P4, Ps, Pt, u0 + mp.mpf(k) / N)
        return -prefactor * bigsum
    elif epsilon ==0:
        return (Kernelsf(Theory, P1, P2, P3, P4,
            Ps, Pt,1)+Kernelsf(Theory, P1, P2, P3, P4,
            Ps, Pt,-1))/2
    else:
        raise ValueError("epsilon must be either 0,1 or -1")

\end{lstlisting}
All the modules that we define in what follows correspond to identities of the fusion kernels that we check. These identities are:
\begin{itemize}
    \item The reflection property of the fusion kernels
    \item The mixing symmetry (see~\eqref{eq:mixingproperty})
    \item The identity derived by Ioannis and Sylvain for c=25 (see~\eqref{eq:RibaultTsiares}).
    \item The pentaton identity (see~\eqref{eq:Pentagoninlittlepgenericrelation}).
    \item The idempotency relation (see~\eqref{eq:idempotencyf}).
\end{itemize}
\begin{lstlisting}[language=Python]
# =========================================
# Identities to be checked for the fusion kernel
# =========================================

def mixingsymmetrieskernelsf(Theory: Theory, P1: complex, P2: complex, P3: complex, P4: complex,
            Ps: complex, Pt: complex, epsilon: int):
    b = Theory.b
    vacuummomentum=(b + 1/b)/2 +0.00001j #for this momentum = P<1,1>, there is a "fake pole" numerically! For this reason, I shifted a bit by giving them a small imaginary part
    return Kernelsf(Theory,P2,P2,P3,P3,vacuummomentum,Pt,0)*Kernelsf(Theory,Pt,P2,Ps,P4,P3,P1,epsilon)/(Kernelsf(Theory,P2,P2,Ps,Ps,vacuummomentum,P1,0))

# This is the identity that Ionnis wrote down in his paper with Sylvain
def RibaultTsiaresspacelikef(Theory: Theory,Ps:complex, Pt:complex, epsilon:int):
    Ps, Pt =mp.mpc(Ps), mp.mpc(Pt)
    m, n = Theory.m, Theory.n
    if (m!=1 or n!=1):
        raise ValueError("The central charge is not 25")
    else:
        if (epsilon==1):
            return 1j*Pt*16**(Pt**2-Ps**2)*mp.exp(2j*mp.pi*Ps*Pt)/Ps
        elif (epsilon==-1):
            return -1j*Pt*16**(Pt**2-Ps**2)*mp.exp(-2j*mp.pi*Ps*Pt)/Ps
        else:
            raise ValueError("epsilon must be either 1 or -1")

# This returns the LHS and the RHS of the Pentagon identity
def Pentagon(Theory: Theory, Param: ParametersPentagon, epsilon: int):
    Po, Pq, P1, P2, P3, P4, P5, Pt, Pu = (
        Param.Po, Param.Pq, Param.P1, Param.P2,
        Param.P3, Param.P4, Param.P5, Param.Pt, Param.Pu
    )

    IntegrandLHS = lambda ps: (
        Kernelsf(Theory, Pu, 1j*(ps+0.0001j), P3, P5, Pq, P4, epsilon)
        * Kernelsf(Theory, Pu, P1, P2, Pq, Po, 1j*(ps+0.0001j), epsilon)
        * Kernelsf(Theory, P1, P2, P3, P4, 1j*(ps+0.0001j), Pt, 0)
    )

    LHS = mp.quad(IntegrandLHS, [-Theory.Lambda, Theory.Lambda])

    RHS = (
        Kernelsf(Theory, Pu, P1, Pt, P5, Po, P4, epsilon)
        * Kernelsf(Theory, Po, P2, P3, P5, Pq, Pt, epsilon)
    )

    return LHS, RHS

def Idempotencyf(Theory: Theory, Param:ParametersIdempotency, epsilon:int):
    Ps, Psprime, P1, P2, P3, P4 = (
        Param.Ps, Param.Psprime, Param.P1, Param.P2,
        Param.P3, Param.P4
    )
    Integrand = lambda p: (
        Kernelsf(Theory, P1, P2, P3, P4, Ps, 1j*(p+0.0001j), epsilon)
        * Kernelsf(Theory, P2, P3, P4, P1, 1j*(p+0.0001j), Psprime, epsilon)
    )

    LHS = mp.quad(Integrand, [-Theory.Lambda, Theory.Lambda])

    return LHS
\end{lstlisting}
We then define the modular kernels. We start by the quantum modular polynomial.
\begin{lstlisting}
# =========================================
# ===================
# Modular kernel
# ===================
# =========================================

# =========================================
# Quantum modular modular polynomial
# =========================================

# quantum modular determinant D
def quantum_determinantm(Theory: Theory, P0:complex, Ps:complex, Pt:complex):
    m, n= Theory.m, Theory.n
    s= Theory.s
    return 1-(smn(Theory, P0/2 - (m+n)/(4*s))/(smn(Theory, Pt)*smn(Theory, Ps)))**2

# roots of the polynomial
def rootsm(Theory: Theory, P0:complex, Ps:complex, Pt:complex):
    m, n= Theory.m, Theory.n
    s= Theory.s
    N = m*n

    prefactor = smn(Theory,Ps)/smn(Theory,P0/2 -Ps -(m+n)/(4*s))

    plusrootfactor = (-1)**N*cmn(Theory,Pt)-1j*smn(Theory,Pt)*mp.sqrt(quantum_determinantm(Theory,P0,Ps,Pt))
    minusrootfactor = (-1)**N*cmn(Theory,Pt)+1j*smn(Theory,Pt)*mp.sqrt(quantum_determinantm(Theory,P0,Ps,Pt))
    
    z_plus= prefactor*plusrootfactor
    z_minus= prefactor*minusrootfactor

    return z_plus, z_minus
\end{lstlisting}
We then construct the modular kernels.
\begin{lstlisting}
    # =========================================
# Modular Kernels
# =========================================

# These are the elementary building blocks of the modular kernels
def CurlyM(Theory: Theory, P0:complex, Ps:complex, Pt:complex, z:complex):
    P0, Ps, Pt, z = mp.mpc(P0), mp.mpc(Ps), mp.mpc(Pt), mp.mpc(z)
    m, n = Theory.m, Theory.n
    s = Theory.s

    term1num = mp.exp(-4j*mp.pi*Ps*z*s)*Gtilde_mn(Theory, z -1+(1/m + 1/n)/4 -(2*Pt+P0)/(2*s))
    term2num = Gtilde_mn(Theory, z -1+(1/m + 1/n)/4 -(-2*Pt+P0)/(2*s))
    Numerator = term1num*term2num

    term1den = Gtilde_mn(Theory, z -(1/m + 1/n)/4 - (2*Pt-P0)/(2*s))
    term2den = Gtilde_mn(Theory, z -(1/m + 1/n)/4 - (-2*Pt-P0)/(2*s))
    Denominator = term1den*term2den

    return Numerator/Denominator
    
def M(Theory: Theory, P0:complex, Ps:complex, Pt:complex):
    P0, Ps, Pt = mp.mpc(P0), mp.mpc(Ps), mp.mpc(Pt)
    m, n = Theory.m, Theory.n
    s = Theory.s
    b = Theory.b
    N = m*n

    Nmn = (-1)**N * (2*mp.pi)**(2*N) * mp.exp(2j*mp.pi*s*Ps)/(s**2)
    Plancherel = -4*mp.sqrt(2)*mp.sin(2*mp.pi*b*Pt)*mp.sin(2*mp.pi*Pt/b)

    term0num=s*(2*mp.pi)**(s*P0-m-n) *Nmn*Plancherel
    term1num=Gmn(Theory,1/m + 1/n +2*Pt/s)
    term2num=Gmn(Theory,(1/m + 1/n)/2 -(P0+2*Ps)/s)
    term3num=Gmn(Theory,1/m + 1/n -2*Pt/s)
    term4num=Gmn(Theory,(1/m + 1/n)/2 -(P0-2*Ps)/s)
    Numerator = term0num*term1num*term2num*term3num*term4num

    term0den=2*Gtilde_mn(Theory,-P0/s)
    term1den=Gmn(Theory,1/m + 1/n + 2*Ps/s)
    term2den=Gmn(Theory,(1/m + 1/n)/2 - (P0+2*Pt)/s)
    term3den=Gmn(Theory,1/m + 1/n - 2*Ps/s)
    term4den=Gmn(Theory,(1/m + 1/n)/2 - (P0-2*Pt)/s)
    Denominator = term0den*term1den*term2den*term3den*term4den
    return Numerator/Denominator

#These are the plus/minus kernels/TV kernels
def KernelsM(Theory: Theory, P0: complex, Ps: complex, Pt: complex, epsilon: int):

    # epsilon =\pm are the elementary kernels while epsilon=0 is the Teschner-Vartanov kernel
    s = Theory.s
    N = Theory.m * Theory.n  # this is s^2 as an integer type

    zplus, zminus = rootsm(Theory, P0, Ps, Pt)

    prefactor = (
        M(Theory, P0, Ps, Pt)
        / (2j * smn(Theory, Pt) * smn(Theory, Ps) * mp.sqrt(quantum_determinantm(Theory, P0, Ps, Pt)))
    )

    bigsum = mp.mpc(0)

    if epsilon == 1:
        u0 = mp.log(zplus) / (2j * mp.pi * s**2)
        for k in range(N):
            bigsum += CurlyM(Theory, P0, Ps, Pt, u0 + mp.mpf(k) / N)
        return prefactor * bigsum

    elif epsilon == -1:
        u0 = mp.log(zminus) / (2j * mp.pi * s**2)
        for k in range(N):
            bigsum += CurlyM(Theory, P0, Ps, Pt, u0 + mp.mpf(k) / N)
        return -prefactor * bigsum
    
    elif epsilon ==0:
        return (KernelsM(Theory, P0, Ps, Pt,1)+KernelsM(Theory, P0, Ps, Pt,-1))/2
    else:
        raise ValueError("epsilon must be either 0,1 or -1")
\end{lstlisting}
All the modules that we define in what follows correspond to identities of the modular kernels. At the moment, these are
\begin{itemize}
    \item The identity limit (see~\eqref{eq:identitylimit}).
    \item The idempotency relation (see~\eqref{eq:idempotencym}).
\end{itemize}
\begin{lstlisting}
#Consistency checks for the modular kernels

def Identitylimit(Theory: Theory, Ps: complex, Pt: complex):
    ident = -(Theory.b+1/Theory.b + 0.00001j)/2 # This is to avoid a fake pole
    return KernelsM(Theory, ident, Ps, Pt, 0), mp.sqrt(2)*mp.cos(4*mp.pi*Ps*Pt)

def Idempotencym(Theory: Theory, P0: complex, P1: complex, P2: complex, epsilon: int ):
    Integrand = lambda p: (
        KernelsM(Theory, P0, P1, 1j*(p+0.00001j), epsilon)*KernelsM(Theory, P0, 1j*(p+0.00001j), P2, epsilon)
    )

    LHS = mp.quad(Integrand, [-Theory.Lambda, Theory.Lambda])

    return LHS
\end{lstlisting}
Here is the battery of tests that we implemented.
\begin{lstlisting}
# =========================================
# Demo tests
# =========================================

if DEMO:
    # --- Demo 1: shift identity youre currently testing ---
    z = mp.mpc(0.6+0.8j)
    z_plus = z + mp.mpf(1) / mp.mpf(theo.n)

    ratio = Gmn(theo, z_plus) / Gmn(theo, z)

    m_mp = mp.mpf(theo.m)
    rhs = mp.power(m_mp, mp.mpf(1)/2 - m_mp * z) \
          * mp.power(2 * mp.pi, (m_mp - mp.mpf(1)) / 2) \
          * mp.gamma(m_mp * z)

    print("=== Demo: Gmn shift check ===")
    print(f"ratio = {ratio}")
    print(f"rhs   = {rhs}")
    print(f"relerr= {abs(ratio-rhs)/max(1,abs(rhs))}\n")

    # --- Demo 2: conjugation & inversion checks for Gtilde_mn ---
    Gv = Gtilde_mn(theo, z)
    print("=== Demo: Gtilde_mn checks ===")
    print(f"conj(Gtilde_mn(z))     = {mp.conj(Gv)}")
    print(f"Gtilde_mn(conj(z))     = {Gtilde_mn(theo, mp.conj(z))}")
    print(f"Gtilde_mn(-z)          = {Gtilde_mn(theo, -z)}")
    print(f"1/Gtilde_mn(z)         = {1/Gv}")
    print()

    # --- Demo 3: alpha at zero point ---
    alpha_demo = alphaf(theo, 0, 0, 0, 0, 0, 0)
    print("=== Demo: alpha(0,0) ===")
    print(f"alpha = {alpha_demo}")

    #---momenta for consistency checks
    P1=0.00001+0.01j
    P2=0.00001+1j
    P3=0.00001+0.03j
    P4=0.00001+0.04j
    Ps=0.00001+0.05j
    Pt=0.00001+0.06j
    P0=0.00001+0.07j
    #--- Demo 4: Check that the roots obtained from the function roots(Theory: Theory,Ps,Pt) are the same that the ones obtain from (-beta \pm sqrt(beta^2-4*alpha*gamma)/2*alpha) ---
    z_plus, z_minus = rootsf(theo, P1, P2, P3, P4, Ps, Pt)
    alpha_val = alphaf(theo, P1, P2, P3, P4, Ps, Pt)
    beta_val = betaf(theo, P1, P2, P3, P4, Ps, Pt)
    gamma_val = gammaf(theo, P1, P2, P3, P4, Ps, Pt)
    z_plus_check = (-beta_val + mp.sqrt(beta_val**2 - 4*alpha_val*gamma_val)) / (2 * alpha_val)
    z_minus_check = (-beta_val - mp.sqrt(beta_val**2 - 4*alpha_val*gamma_val)) / (2 * alpha_val)
    print("=== Demo: roots consistency check ===")
    print(f"z_plus from roots() = {z_plus}")
    print(f"z_plus from formula   = {z_plus_check}")
    print(f"z_minus from roots() = {z_minus}")
    print(f"z_minus from formula   = {z_minus_check}")

    #--- Demo 5: Check the symmetry of CurlyF under the transformation P_s\rightarrow P_3,P_3\rightarrow P_s,P_t\rightarrow P_1,P_1\rightarrow P_t ---
    CurlyF_original = CurlyF(theo, P1, P2, P3, P4, Ps, Pt, z)
    CurlyF_transformed = CurlyF(theo, Pt, P2, Ps, P4, P3, P1, z)
    print("=== Demo: CurlyF symmetry check ===")
    print(f"CurlyF original     = {CurlyF_original  }")
    print(f"CurlyF transformed  = {CurlyF_transformed}")

    # --- Demo 6: Check the single-valuedness of the sum in the Kernels function ---
    z_plus, z_minus = rootsf(theo, P1, P2, P3, P4, Ps, Pt)

    N = theo.m * theo.n
    s = theo.s
    ell = 3

    res1 = mp.mpc(0)
    res2 = mp.mpc(0)

    u0 = mp.log(z_plus) / (2j * mp.pi * s**2)

    for k in range(N):
        z1 = u0 + mp.mpf(k) / N
        z2 = u0 + mp.mpf(k + ell) / N
        res1 += CurlyF(theo, P1, P2, P3, P4, Ps, Pt, z1)
        res2 += CurlyF(theo, P1, P2, P3, P4, Ps, Pt, z2)

    print("=== Demo: CurlyF single-valuedness check ===")
    print(f"sum unshifted = {res1}")
    print(f"sum shifted   = {res2}")

    #--- Demo7: Check the reflection property of the spacelike kernels.
    Fplus=Kernelsf(theo,P1,P2,P3,P4,Ps,Pt,1)
    Fminus=Kernelsf(theo,P1,P2,P3,P4,Ps,Pt,-1)
    Fplusreflected=Kernelsf(theo,P1,P2,P3,P4,-Ps,Pt,1)
    Fminusreflected=Kernelsf(theo,P1,P2,P3,P4,-Ps,Pt,-1)
    print("=== Demo: Reflection property of the plus/minus kernels ===")
    print(f"Fplus = {Fplus}")
    print(f"Fminus   = {Fminus}")
    print(f"Fplusreflected = {Fplusreflected}")
    print(f"Fminusreflected   = {Fminusreflected}")

    #---Demo8: Check the reflection property of the Teschner-Vartanov kernel
    print("=== Demo: Reflection property of Teschner-Vartanov kernel ===")
    print(f"FTeschnerVartanov = {Kernelsf(theo,P1,P2,P3,P4,Ps,Pt,0)}")
    print(f"FTeschnerVartanovreflected = {Kernelsf(theo,P1,P2,P3,P4,-Ps,Pt,0)}")

    #---Demo9: Check the mixing symmetries of the different kernels
    Fplusmixed=mixingsymmetrieskernelsf(theo,P1,P2,P3,P4,Ps,Pt,1)
    Fminusmixed=mixingsymmetrieskernelsf(theo,P1,P2,P3,P4,Ps,Pt,-1)
    print("=== Demo: mixing property of the kernels ===")
    print(f"Fplus = {Fplus}")
    print(f"Fplusmixed = {Fplusmixed}")
    print(f"Fminus = {Fminus}")
    print(f"Fminusmixed = {Fminusmixed}")

    #---Demo10: Check Ioannis/Sylvain relation for the spacelike kernels
    Fplusatonefourth=Kernelsf(theo,0.25,0.25,0.25,0.25,Ps,Pt,1)
    Fminusatonefourth=Kernelsf(theo,0.25,0.25,0.25,0.25,Ps,Pt,-1)
    Fplusfromrelation=RibaultTsiaresspacelikef(theo,Ps,Pt,1)
    Fminusfromrelation=RibaultTsiaresspacelikef(theo,Ps,Pt,-1)
    print("=== Demo: Ioannis/Sylvain relation for the spacelike kernels at c=25 ===")
    print(f"Fplusatonefourth = {Fplusatonefourth}")
    print(f"Fplusfromrelation = {Fplusfromrelation}")
    print(f"Fminusatonefourth = {Fminusatonefourth}")
    print(f"Fminusfromrelation = {Fminusfromrelation}")

    # ---Demo11: Check pentagon for the TV kernel
    print("=== Demo: Testing the pentagon for the TeschnerVartanov kernel ===")
    LHS, RHS = Pentagon(theo, parPenta, 0)
    print(f"LHSPentagonplus = {LHS}")
    print(f"RHSPentagonplus = {RHS}")

    # ---Demo12: Check Idempotency for the TV kernel
    print("=== Demo: Testing the Idempotency for the TeschnerVartanov kernel ===")
    LHS= Idempotencyf(theo, parIdem, 0)
    print(f"Idempotency = {LHS}")

    # --- Demo 13: Check the single-valuedness of the sum in the modular Kernels function ---
    z_plus, z_minus = rootsm(theo, P0, Ps, Pt)

    N = theo.m * theo.n
    s = theo.s
    ell = 3

    res1 = mp.mpc(0)
    res2 = mp.mpc(0)

    u0 = mp.log(z_plus) / (2j * mp.pi * s**2)

    for k in range(N):
        z1 = u0 + mp.mpf(k) / N
        z2 = u0 + mp.mpf(k + ell) / N
        res1 += CurlyM(theo, P0, Ps, Pt, z1)
        res2 += CurlyM(theo, P0, Ps, Pt, z2)

    print("=== Demo: CurlyM single-valuedness check ===")
    print(f"sum unshifted = {res1}")
    print(f"sum shifted   = {res2}")

    # #--- Demo 14: Check the Identity limit of the Modular kernel
    LHS, RHS = Identitylimit(theo, Ps, Pt)
    print("=== Demo: Identity limit check ===")
    print(f"Kernel at identity = {LHS}")
    print(f"Value of the cos   = {RHS}")

    # #--- Demo 15: Check reflection properties under reflection of the external momenta
    print("=== Demo: Reflection property of Modular kernels ===")
    print(f"Mplus = {KernelsM(theo,P0,Ps,Pt,1)}")
    print(f"Mminus = {KernelsM(theo,P0,Ps,Pt,-1)}")
    print(f"Mplusreflectedps = {KernelsM(theo,P0,-Ps,Pt,1)}")
    print(f"Mminusreflectedps = {KernelsM(theo,P0,-Ps,Pt,-1)}")
    print(f"Mplusreflectedpt = {KernelsM(theo,P0,Ps,-Pt,1)}")
    print(f"Mminusreflectedpt = {KernelsM(theo,P0,Ps,-Pt,-1)}")
    print(f"MTV = {KernelsM(theo,P0,Ps,Pt,0)}")
    print(f"MTVreflectedps = {KernelsM(theo,P0,-Ps,Pt,0)}")
    print(f"MTVreflectedpt = {KernelsM(theo,P0,Ps,-Pt,0)}")

\end{lstlisting}
Here I display the results of the tests.
\begin{lstlisting}
=== Demo I: Consistency checks for the special functions ===

=== Demo: Gmn shift check ===
ratio = (0.49438518268587436096 - 0.48170611659036531661j)
rhs   = (0.49438518268587436096 - 0.48170611659036531661j)
=== Demo: Gtilde_mn checks ===
conj(Gtilde_mn(z))     = (0.48298133429787586343 - 0.46023939521263494923j)
Gtilde_mn(conj(z))     = (0.48298133429787586343 - 0.46023939521263494923j)
Gtilde_mn(-z)          = (1.0851287514495764127 - 1.0340337500226893283j)
1/Gtilde_mn(z)         = (1.0851287514495764127 - 1.0340337500226893283j)

=== Demo II: Consistency checks for the quantum modular fusion polynomial ===

=== Demo: alpha(0,0) ===
alpha = (0.0 + 4.1924144845297254444e-21j)
=== Demo: roots consistency check ===
z_plus from roots() = (-1.2611678633146055076 + 0.000049540218341473836137j)
z_plus from formula   = (-1.2611678633146055075 + 0.000049540218341473822723j)
z_minus from roots() = (-0.96838174127724644965 + 0.000062525639169306944851j)
z_minus from formula   = (-0.96838174127724644976 + 0.000062525639169306958265j)

=== Demo III: Consistency checks for the spacelike fusion kernels ===

=== Demo: CurlyF symmetry check ===
CurlyF original     = (1.0004160450050665359 + 0.02442251129663564322j)
CurlyF transformed  = (1.0004160450050665359 + 0.02442251129663564322j)
=== Demo: CurlyF single-valuedness check ===
sum unshifted = (0.9999938049400514083 - 0.0028236456207856117867j)
sum shifted   = (0.99999380494005140954 - 0.0028236456207856147357j)
=== Demo: Reflection property of the plus/minus kernels ===
Fplus = (0.0070379226124723145446 + 6.7370768266815472853j)
Fminus   = (0.0081042804008412503518 - 6.7370818422901804711j)
Fplusreflected = (0.0081042804008412502499 - 6.7370818422901804711j)
Fminusreflected   = (0.0070379226124723090828 + 6.7370768266815472775j)
=== Demo: Reflection property of Teschner-Vartanov kernel ===
FTeschnerVartanov = (0.0075711015066567824482 - 2.5078043165928696296e-6j)
FTeschnerVartanovreflected = (0.0075711015066567796664 - 2.5078043165968235794e-6j)
=== Demo: mixing property of the kernels ===
Fplus = (0.0070379226124723145446 + 6.7370768266815472853j)
Fplusmixed = (0.0081054953745872160188 - 6.7370818411835346656j)
Fminus = (0.0081042804008412503518 - 6.7370818422901804711j)
Fminusmixed = (0.0070367076404285165413 + 6.737076828305664042j)
=== Demo: Ioannis/Sylvain relation for the spacelike kernels at c=25 ===
Fplusatonefourth = (0.022508560047440773346 + 1.1961257192910219677j)
Fplusfromrelation = (0.022508560047440773356 + 1.1961257192910219677j)
Fminusatonefourth = (0.02258994051224826017 - 1.1961407251585321349j)
Fminusfromrelation = (0.02258994051224826016 - 1.1961407251585321348j)
=== Demo: Testing the pentagon for the TeschnerVartanov kernel ===
LHSPentagonplus = (0.040478458528855747172 + 0.000094724396361210183938j)
RHSPentagonplus = (0.040480286896388033099 + 0.000094758445189765142584j)
=== Demo: Testing the Idempotency for the TeschnerVartanov kernel ===
Idempotency = (-1.3517102240608496172e-8 - 9.1919226111705190442e-8j)

=== Demo IV: Consistency checks for the spacelike modular kernels ===

=== Demo: CurlyM single-valuedness check ===
sum unshifted = (1.3609325834244597261 - 0.14924114543484447604j)
sum shifted   = (1.3609325834244597259 - 0.14924114543484447632j)
=== Demo: Identity limit check ===
Kernel at identity = (1.4132087274245876136 - 6.9794754220596457469e-6j)
Value of the cos   = (1.4132087251469752577 + 7.3679348136474341459e-7j)
=== Demo: Reflection property of Modular kernels ===
Mplus = (0.063138621188207340149 - 0.82235300568100736394j)
Mminus = (0.063378389776097701905 + 0.8223112349827522751j)
Mplusreflectedps = (0.063378389776097701907 + 0.8223112349827522751j)
Mminusreflectedps = (0.0631386211882073401 - 0.82235300568100736392j)
Mplusreflectedpt = (0.063378389776097701906 + 0.8223112349827522751j)
Mminusreflectedpt = (0.063138621188207340149 - 0.82235300568100736394j)
MTV = (0.063258505482152521027 - 0.000020885349127544417484j)
MTVreflectedps = (0.063258505482152521003 - 0.00002088534912754440605j)
MTVreflectedpt = (0.063258505482152521027 - 0.000020885349127544417484j)

=== Demo V: important identities involving the spacelike modular and fusion kernels ===

=== Demo: Non rational Verlinde formula ===
LHSVerlindeformula = (0.0018437816707970328695 - 1.2143930559352561683e-6j)
RHSVerlindeformula = (0.0018361623980578999904 - 1.2095319899761086453e-6j)
\end{lstlisting}


As you can see, I did a plot of the quantum modular determinant in the complex. As you can observe, a very strong pattern is observed! Exercise: try to understand this very interesting pattern...
\begin{figure}[t!]
    \includegraphics[width=0.8\textwidth]{Quantummodulardeterminant.png}
    \caption{Phase plot of the quantum determinant in the $P_s$ complex plane. The parameters are $P_1=0.2i$, $P_2=0.3i$, $P_3=0.4i$, $P_4=0.5i$, $P_t=0.7i$ and $m=n=1$.}
\end{figure}

\end{document}